\section{Discussion and Limitations}
\label{sec:discussion}

Our analysis (\S\ref{sec:analysis}) leaves a clear practical implication: real headroom on this
benchmark requires structurally stronger \emph{visual} models (end-to-end fine-tuning that
preserves the prior, multi-crop test-time augmentation, test-time training, synthetic mosaic
training distributions) rather than post hoc reweighting with auxiliary priors. The remainder of
this section is candid about what our results do and do not establish.

\paragraph{What we did not cleanly ablate.} Our strongest model, i002, differs from the earlier
010 fine-tune along several axes at once: the training data ($1.41$M $\to$ $2.65$M images, a
different manifest), the schedule (five $\to$ ten epochs and a larger effective batch), the head
topology (a shared MLP trunk $\to$ independent per-head MLPs), and the head count (five $\to$
three taxonomic levels). No run isolates any one of these factors, so we report i002 as the
stronger configuration overall and make \emph{no} causal claim that the per-head head design
caused the gain. Nor do we ablate the taxonomy supervision itself: we never trained a
species-only control with the auxiliary heads removed, so the auxiliary losses are reported as a
deliberate regulariser, not as a measured improvement. The one head-adjacent factor we did
isolate, the number of unfrozen blocks (holding the 010 head fixed), shows a sharp optimum at
four blocks (\S\ref{sec:blocks_ablation}), and that is the surface the final anchor uses. The
single near-controlled counterfactual we have is the i003 cap (\S\ref{sec:longtail}), where only
the manifest changed and the score dropped.

\paragraph{Public and private leaderboard instability.} The public and private splits do not rank
our configurations the same way. Our best public submission ($0.41826$) is not our best private
one: the lightly logit-adjusted single-resolution run is the strongest on the private split
($0.40600$), and an i002+010 cross-checkpoint ensemble is near-best private while weak public.
The three pivot deltas ($-0.005$ to $-0.036$) are mostly within or near the ${\sim}0.002$
run-to-run noise floor we characterise in \S\ref{sec:eval_protocol}, so apart from cRT they should
be read as ``no usable effect'' rather than as measured decreases. We therefore carry both
columns throughout the paper and avoid implying that the best public run is best overall.

\paragraph{Limitations of the representation analysis.} The head-representation study
(\S\ref{sec:head_repr}) is descriptive geometry, not a controlled experiment. It uses a single
checkpoint per design, a single $4{,}000$-image sample per experiment, and two different
validation splits (only $62$ of $4{,}000$ sampled images overlap), so we compare each
representation only against its own encoder and never cross-experiment. It characterises what the
shared and per-head designs do to the feature space; it does not isolate the effect of head
topology on Macro F1. More broadly, the full development trace in
Appendix~\ref{app:experiment_summary} is exactly that, a trace rather than a controlled
ablation table, and most of its rows change several factors at once.

\paragraph{Future work.} Several directions follow directly from these findings. The first is to
pull on the data axis rather than the capacity axis: a synthetic mosaic training regime that
composes single-plant images into multi-species scenes with controlled overlap would let the
network see the quadrat distribution during training, which neither the PlantCLEF~2024 corpus nor
iNaturalist habitat shots do at scale. The second is self-distillation against our own production
checkpoint, using its predictions on the unlabelled LUCAS pseudo-quadrat corpus as soft targets
filtered by confidence and entropy. The third is genuinely orthogonal evidence that is not in the
image at all, in particular a location-conditioned bioclimatic envelope prior built from
training-image coordinates against WorldClim variables at each test site, the one orthogonal axis
that visual features cannot recover from the image alone (this experiment is blocked here only by
the absence of test-site coordinates). Finally, per-class threshold calibration on a held-out
multi-species validation set would target Macro F1 more directly than the global probability
threshold we currently use, and could absorb residual species-frequency imbalance that our
manifest leaves behind.
